We begin by defining what is an Algebra. This is the main structure we will be using throughout this paper. To give you an intuition first, an Algebra is basically a set with a structure of both a ring an a vector space.

\section{Algebras over a Field}


\begin{definition}[k-Algebra]
Let $k$ be any field. We say that a $k-$vector space $A$ is a $k$-algebra if $A$ is equipped with an associative bilinear operation (product)
\[
\times : A \times A \to A
\]
satisfying the following axioms of bilinearity.
For every $x,y,z \in A$ and $a,b \in k$:
\begin{enumerate}
    \item 
    
$    z \times (x+y) = z \times x + z\times y
$    
    \item
    
$    (x+y)\times z = x\times z + y \times z
$    
    \item
    
$    (a \cdot x)\times (b \cdot y) = ab \cdot (x\times y)
$    
    
\end{enumerate}
where $\cdot $ is the vector space scalar multiplication.
\end{definition}

Another way to define a $k$-algebra is to first define it as a ring equipped with scalar multiplication operation by embedding the field $k$ inside $A$.

\begin{definition}[k-Algerba]
    Let $k$ be a field. A $k-$algebra $A$ is a ring with an injective ring homomorphism $k \to A$, where for any $\lambda \in k$ and $a \in A$ we have that 
    \[
    \lambda \cdot a = a \cdot \lambda
    \]
i.e. multiplication by elements in the image of $k$ in $A$ is commutative.
\end{definition}

\large \textbf{Remarks:-}\\
\begin{enumerate}
    \item Observe that we can define the vector space structure on $A$ in the definition above by defining the scalar product $\lambda a $ to be the product of the image of $\lambda$ in $A$ times $a$ to get $\lambda \cdot a$ for all $\lambda \in k$ and $a \in A$. That is
    \[\lambda a = \lambda \cdot a\]
    \item An algebra product operation might not be commutative. If so we just say a commutative algebra.
\item We may or may not include $1$ as a unity in our definition of an algebra but here we do. And the operation is associative.

\end{enumerate}



\begin{definition}[k-Algebra Homomorphism]
    Let $A,B$ be $k-$algebras. A function $\varphi: A \to B$ is a $k-$algebra homomorphism if it is a linear map of vector spaces and for any $a \in A$ and $b \in B$ we have that 
    \[\varphi(ab) = \varphi(a) \varphi(b)\]
    
\end{definition}

\section{Direct Sum and Direct Product}

Now, we get to differentiate between two important notions in linear algebra: Internal Direct Sum and External Direct Product. Sometimes for simplicity we just say Direct Sum and Direct Product.

\begin{definition}[Internal Direct Sum]
Let $V$ be a vector space. We say $V$ is an internal direct sum of two subspaces $U,W \subseteq V$ iff every vector $v\in V$ could be written uniquely as a sum of two vectors $u,w$
where $u\in U$, and $w\in W$.

That is
\[
V=U+W.
\]
there are unique vectors $u\in U$ and $w\in W$ s.t.
\[
v=u+w.
\]
\end{definition}

The above definition tells us exactly when a sum of two subspaces $U+W$ of a given vector space $V$ is a direct sum, and in this case we call $V$ an internal direct sum since the two spaces $U$ and $W$ are subspaces. Next, we define an External Direct Product of two separate given spaces.

\begin{definition}[External Direct Product]
Let $K,L$ be two vector spaces. We define the \emph{External Direct Product} of $K$ and $L$ to be the vector space
\[
K\times L:=\{(k,l)\mid k\in K,\ l\in L\}
\]
with operations of component wise addition and scalar multiplication.
\end{definition}


The next theorem provides a complete characterization for when an External Direct Product of two subspaces and their sum are exactly the same, giving also an equivalent condition for when a sum is an internal direct sum.

\begin{theorem}
Let $U,W$ be subspaces of a vector space $V$. \textbf{TFAE}

\begin{enumerate}
    \item The map
    \[
    \pi:U\times W\to U+W
    \]
    given by
    \[
    (u,w)\mapsto u+w
    \]
    is an isomorphism.

    \item $ U\cap W=\{0\}.$
    
    \item Any vector $x\in U+W$ could be written uniquely in the form
    \[
    u+w
    \]
    with $u\in U$ and $w\in W$.
\end{enumerate}
\end{theorem}

\begin{proof}
$(1)\Rightarrow(2)$ Suppose for the sake of contradiction that there exist a non-zero vector
\[
v\in U\cap W.
\]
So $v\in U$ and $v\in W$. Hence via the isomorphism $\pi$, the element $(v,0) \in U \times W$ maps to $v+0 = v$ and $(v,0) \in U \times W$ maps to $0+v = v$ violating injectivity of $\pi$.

$(2)\Rightarrow(3)$ Suppose that there is $x\in U+W$ such that $x$ has two different representations
\[
x=a+b=c+d.
\]
where $a,c \in U$ and $b,d \in W$ 
Then
\[
a+b=c+d
\]
\[
\Rightarrow a-c=d-b.
\]

Since $U,W$ are subspaces that is closed under addition we get that 
$a-c\in U$
and
$b-d\in W.$ Since $a-c=d-b$, therefore,
\[
a-c\in U\cap W.\] 
which implies that $a-c= d-b= 0$. Hence, $a=c$ and $b=d$ as desired.


$(3)\Rightarrow(1)$ Since any vector $ x \in U +W$ could be written as $u+w$, where $u \in U$ and $w \in W$ this implies that $\pi$ is surjective as a preimage for any element $x = u+w$ is $(u,w)$. By uniqueness of this representation of elements in $U+W$ we get that the preimage of an element $x = u+w$ is unique.
\end{proof}

So, if we have a vector space that is a sum of two subspaces
\[
U+W=V
\]
and satisfies one of the conditions above we say $V$ is an internal direct sum of $U,W$ and we write
\[
V=U\oplus W.
\]
Moreover, we can see that the two notions of internal direct sum and external direct product coincide when $K,L$ are subspaces of some vector space $V$ and $V$ is an external direct sum of $K$ and $L$ as proved the direct sum \& product are isomorphic.


\begin{definition}[Direct sum of Algebras]
    Let $A,B$ be two $k-$algebras. We define the direct sum $A \oplus B$ to be the $k-$algebra whose elements are pairs $(a,b)$ where $a \in A$ and $b \in B$, with component wise operations.
\end{definition}
% Now, another linear algebraic notion we need here is the Jordan form.

% \begin{definition}[$k$-Algebra]
% \label{def:k_algebra}
% Let $k$ be any field. We say a set $A$ is a \emph{$k$-algebra} if $A$ is a $k$-vector space equipped with a bilinear operation
% \[
% \cdot : A \times A \longrightarrow A
% \]
% satisfying the following axioms for all $x, y, z \in A$ and all $a, b \in k$:
% \begin{enumerate}
%     \item $z \cdot (x + y) = z \cdot x + z \cdot y$ \quad \text{(left distributivity)},
%     \item $(x + y) \cdot z = x \cdot z + y \cdot z$ \quad \text{(right distributivity)},
%     \item $(ax) \cdot (by) = ab(x \cdot y)$ \quad \text{(bilinearity / compatibility with scalar multiplication)}.
% \end{enumerate}
% \end{definition}

% \subsection{Alternative Formulation via Rings}
% Another equivalent way to define a $k$-algebra is to first define it as a ring $(A, +, \cdot)$ equipped with a scalar multiplication operation $k \times A \to A$ such that $A$ is a $k$-vector space and scalar multiplication is compatible with the ring product:
% \[
% c(x \cdot y) = (cx) \cdot y = x \cdot (cy) \quad \text{for all } c \in k, \; x, y \in A.
% \]
% Equivalently, a $k$-algebra is a ring $A$ with a ring homomorphism $\eta: k \longrightarrow Z(A)$ from $k$ into the center of $A$.


% \text{Remarks:-}\\
% An algebra is not necessarily commutative under multiplication. When the multiplication satisfies $xy = yx$ for all $x, y \in A$, we refer to it as a \emph{commutative algebra}.

% In general algebraic literature, one may or may not include a multiplicative identity $1$ (a unity) in the basic definition. \textbf{In these notes, we explicitly assume the presence of a unity $1_A$}, and we assume the multiplication is associative:
% \[
% (x \cdot y) \cdot z = x \cdot (y \cdot z) \quad \text{for all } x, y, z \in A.
% \]


% \section{Direct Sums and Direct Products in Linear Algebra}

% Now, we differentiate between two foundational notions in linear algebra: \emph{internal direct sums} and \emph{external direct products}.

% \begin{definition}[Direct Sum of Subspaces]
% \label{def:direct_sum}
% Let $V$ be a vector space over $k$. We say $V$ is a \emph{direct sum} of two subspaces $U, W \subseteq V$ if and only if every vector $v \in V$ can be written uniquely as a sum of two vectors $u$ and $w$, where $u \in U$ and $w \in W$.
% That is,
% \[
% \forall v \in V, \quad \exists!\, u \in U, \; w \in W \quad \text{such that} \quad v = u + w.
% \]
% \end{definition}

% The above definition characterizes precisely when the sum of two subspaces $U + W$ is considered a direct sum. The following theorem provides a complete algebraic characterization.

% \begin{theorem}[Characterization of Direct Sums]
% \label{thm:direct_sum_tfae}
% Let $U$ and $W$ be subspaces of a vector space $V$. The following statements are equivalent (\textbf{TFAE}):
% \begin{enumerate}
%     \item The linear map
%     \[
%     \pi : U \times W \longrightarrow U + W, \qquad (u, w) \longmapsto u + w
%     \]
%     is an isomorphism of vector spaces.
%     \item $U \cap W = \{0\}$.
%     \item Any vector $x \in U + W$ can be written uniquely in the form $u + w$ with $u \in U$ and $w \in W$.
% \end{enumerate}
% \end{theorem}

% \begin{proof}
% We prove the chain of implications $(1) \implies (2) \implies (3) \implies (1)$.

% \medskip
% \noindent \textbf{(1) $\implies$ (2):} \\
% Suppose for the sake of contradiction that there exists a non-zero vector $v \in U \cap W$.
% Since $v \in U$ and $v \in W$, we may construct two distinct elements in the Cartesian product $U \times W$:
% \[
% (v, 0) \in U \times W \quad \text{and} \quad (0, v) \in U \times W.
% \]
% Evaluating the map $\pi$ on these elements gives:
% \[
% \pi(v, 0) = v + 0 = v,
% \]
% and
% \[
% \pi(0, v) = 0 + v = v.
% \]
% Since $v \neq 0$, we have $(v, 0) \neq (0, v)$, yet $\pi(v, 0) = \pi(0, v)$. This directly violates the injectivity of $\pi$, contradicting our hypothesis that $\pi$ is an isomorphism. Therefore, no such non-zero vector $v$ can exist, and we conclude:
% \[
% U \cap W = \{0\}.
% \]

% \medskip
% \noindent \textbf{(2) $\implies$ (3):} \\
% Suppose for the sake of contradiction that there is a vector $x \in U + W$ with two distinct representations:
% \[
% x = a + b = c + d, \quad \text{where } a, c \in U \text{ and } b, d \in W.
% \]
% Rearranging terms by subtracting $c$ and $b$ from both sides gives:
% \[
% a - c = d - b.
% \]
% Because $U$ is a subspace and $a, c \in U$, the vector $a - c \in U$. Likewise, because $W$ is a subspace and $b, d \in W$, the vector $d - b \in W$ (note: this is equivalent to $b - d \in W$).
% Since the two sides are equal, this common vector belongs to the intersection:
% \[
% a - c = d - b \in U \cap W.
% \]
% By hypothesis (2), $U \cap W = \{0\}$. Thus:
% \[
% a - c = 0 \implies a = c, \qquad \text{and} \qquad d - b = 0 \implies b = d.
% \]
% Hence, the two representations are identical, proving uniqueness.

% \medskip
% \noindent \textbf{(3) $\implies$ (1):} \\
% The map $\pi : U \times W \to U + W$ is linear by the properties of vector addition. It is surjective by the definition of the sum of subspaces $U + W = \{ u + w \mid u \in U, w \in W \}$.
% If $\pi(u, w) = 0$, then $u + w = 0 + 0$. By the uniqueness of the representation of the zero vector guaranteed by (3), we must have $u = 0$ and $w = 0$. Hence, $\ker(\pi) = \{(0, 0)\}$, proving that $\pi$ is injective. Being a bijective linear transformation, $\pi$ is an isomorphism.
% \end{proof}

% \begin{definition}[Direct Sum Notation]
% If a vector space $V$ is the sum of two subspaces $U + W = V$ and satisfies any of the equivalent conditions of Theorem~\ref{thm:direct_sum_tfae}, we say $V$ is the \emph{internal direct sum} of $U$ and $W$, and we write:
% \[
% V = U \oplus W.
% \]
% \end{definition}

% \section{External Direct Product of Vector Spaces}

% Now, the other complementary linear algebraic notion we introduce is the \emph{direct product}.

% \begin{definition}[External Direct Product]
% \label{def:direct_product}
% Let $K$ and $L$ be two vector spaces over the field $k$. We define the vector space
% \[
% K \times L := \{ (k, l) \mid k \in K, \; l \in L \}
% \]
% equipped with componentwise operations:
% \begin{align*}
% (k_1, l_1) + (k_2, l_2) &:= (k_1 + k_2, \; l_1 + l_2), \\
% \alpha (k, l) &:= (\alpha k, \; \alpha l) \quad \text{for all } \alpha \in k, \; k \in K, \; l \in L,
% \end{align*}
% to be the \emph{direct product} of $K$ and $L$.
% \end{definition}

% \begin{lemma}[Equivalence of Internal and External Notions]
% When $K$ and $L$ are subspaces of an ambient vector space $V$ satisfying $K \cap L = \{0\}$ and $K + L = V$, the internal direct sum $K \oplus L$ and the external direct product $K \times L$ are canonically isomorphic:
% \[
% K \times L \xrightarrow{\ \sim\ } K \oplus L, \qquad (k, l) \longmapsto k + l.
% \]
% \end{lemma}

% % \section{Transition to Jordan Canonical Form}

% % Now, another fundamental linear algebraic tool required for representation theory and module decomposition is the \textbf{Jordan canonical form}.

% % Role in Representation Theory:-\\
% % The Jordan canonical form allows the classification of linear endomorphisms $T \in \operatorname{End}(V)$ up to similarity over algebraically closed fields. In the context of module theory:
% % \begin{enumerate}
% %     \item It provides the canonical decomposition of generalized eigenspaces:
% %     \[
% %     V = \bigoplus_{\lambda \in \operatorname{Spec}(T)} \ker(T - \lambda I)^{\dim V}.
% %     \]
% %     \item It decomposes an operator into its semisimple and nilpotent components ($T = D + N$ with $DN = ND$), identifying non-semisimple representations where Jordan blocks of size $> 1$ manifest non-trivial extensions of modules.
% % \end{enumerate}




